\chapter{Synchronisation NTP}

\section{Importance du PDC}

Le PDC Emulator doit être le serveur NTP fiable pour le domaine. Sans synchronisation correcte :

\begin{itemize}
    \item Kerberos échoue (fenêtres de validité des tickets)
    \item Certificats sont rejetés (validation temporelle)
    \item Réplication AD échoue
    \item Événements de sécurité mal horodatés
\end{itemize}

\section{Configuration sur srv-dc01 (PDC)}

Configurer srv-dc01 comme serveur NTP Stratum 2 fiable.

\subsection{Ajouter les sources NTP}

\begin{lstlisting}[style=powershell]
# Configurer les pools NTP publics
w32tm /config /manualpeerlist:"0.pool.ntp.org,0x9 1.pool.ntp.org,0x9" `
    /syncfromflags:manual /reliable:YES /update
\end{lstlisting}

\subsection{Redémarrer le service}

\begin{lstlisting}[style=powershell]
net stop w32time
net start w32time
\end{lstlisting}

\subsection{Forcer la synchronisation}

\begin{lstlisting}[style=powershell]
w32tm /resync /force
\end{lstlisting}

\section{Vérification}

\subsection{Source NTP actuelle}

\begin{lstlisting}[style=powershell]
w32tm /query /source

# Résultat attendu : 0.pool.ntp.org ou 1.pool.ntp.org
\end{lstlisting}

\subsection{État complet}

\begin{lstlisting}[style=powershell]
w32tm /query /status

# Exemple de sortie :
# Leap Indicator: 0 (no warning)
# Stratum: 3
# Precision: -20
# Root Delay: 0.0127440s
# Root Dispersion: 0.1154056s
# ReferenceId: 0xC0A82801 (DC01 — stratum 2)
# Last Successful Sync Time: 12/23/2026 14:53:42
# Source: 0.pool.ntp.org
# Poll Interval: 10 (1024s)
\end{lstlisting}

\subsection{Lister les peers}

\begin{lstlisting}[style=powershell]
w32tm /query /peers

# Affiche les serveurs NTP configurés et leur état
\end{lstlisting}

\section{Configuration des clients}

Les clients (DC secondaire, serveurs, postes) synchronisent automatiquement sur le PDC via Kerberos.

\textbf{Aucune configuration supplémentaire requise}. Windows utilise le protocole de sécurité NT (sécurisé) pour synchroniser avec le PDC.

\section{Troubleshooting NTP}

\subsection{Synchronisation échouée}

\begin{lstlisting}[style=powershell]
# Vérifier la connectivité vers les serveurs NTP
Test-NetConnection -ComputerName 0.pool.ntp.org -Port 123

# Vérifier l'état du service W32Time
Get-Service W32Time | Select-Object Status

# Voir les logs
Get-EventLog -LogName System | Where-Object Source -eq "W32Time" | Tail -20
\end{lstlisting}

\subsection{Écart de temps important}

\begin{lstlisting}[style=powershell]
# Vérifier l'heure système
Get-Date

# Vérifier l'écart par rapport à un serveur
w32tm /stripchart /computer:0.pool.ntp.org /samples:5

# Si écart > 15 minutes, impossible à synchroniser
# → Fixer manuellement l'heure système
Set-Date -Date "2026-05-23 14:30:00"
\end{lstlisting}

\subsection{Timezone incorrect}

\begin{lstlisting}[style=powershell]
# Afficher la timezone actuelle
Get-TimeZone

# Changer la timezone (exemple : Maroc)
Set-TimeZone -Id "Morocco Standard Time"

# Redémarrer W32Time
Restart-Service W32Time
\end{lstlisting}

\section{Checklist NTP}

\begin{itemize}
    \item[\checkmark] Identifier le PDC (Emulator FSMO)
    \item[\checkmark] Configurer sources NTP sur PDC
    \item[\checkmark] Redémarrer W32Time
    \item[\checkmark] Vérifier source avec \texttt{w32tm /query /source}
    \item[\checkmark] Vérifier status complet
    \item[\checkmark] Attendre 10 minutes de stabilisation
    \item[\checkmark] Tester depuis un client
\end{itemize}
